% /Users/pikachu/books/payments-book/src/parts/part5/ch29-issuer-acquirer.tex
\chapter{Эмитент и эквайер: две стороны одной транзакции}
\label{ch:issuer-acquirer}

\begin{kaobox}[title=Что вы узнаете из этой главы]
  Одна и та же транзакция выглядит совершенно по-разному с двух
  сторон: для эмитента это вопрос \enquote{платить или нет},
  для эквайера — \enquote{как обслужить мерчанта}. Какие системы
  стоят за каждой из сторон: CMS, authorization host, ACS —
  у эмитента; TMS, onboarding, merchant monitoring —
  у эквайера; и почему модель Payment Facilitator изменила
  правила игры для онлайн-платформ.
\end{kaobox}

Одна и та же транзакция задаёт двум банкам два разных набора
вопросов. Эмитент спрашивает: можно ли доверить этой операции
деньги и риск? Эквайер спрашивает: как подключить, обслужить
и контролировать мерчанта так, чтобы платёжный поток не развалился
на масштабе. Эту главу стоит читать как симметричный разбор двух
сторон одной записи.

\section{Эмитент}
\label{sec:ia-issuer}

Эмитент (issuer) — банк, выпустивший карту и несущий кредитный
риск по каждой транзакции: когда держатель карты совершает
покупку, именно эмитент гарантирует оплату эквайеру через
карточную сеть, а затем взыскивает средства с держателя —
списывая с дебетового счёта или увеличивая задолженность
по кредитному. За этой простой формулировкой стоит целый комплекс
систем, каждая из которых отвечает за свой аспект жизненного
цикла карты и транзакции.

\subsection*{Card Management System}

Card Management System (CMS) — центральная система эмитента, управляющая жизненным циклом карты от момента заказа до
блокировки и перевыпуска. CMS хранит данные карты (PAN, срок действия, CVV, статус, лимиты, привязка к счёту),
генерирует заявки на персонализацию (запись данных на чип и магнитную полосу) и управляет состояниями: активна,
заблокирована временно, заблокирована окончательно, перевыпущена, истекла.

Персонализация — запись данных на физическую карту. Она включает
загрузку EMV-апплета, ключей эмитента (IMK), параметров рисков
(floor limit, офлайн-лимиты) и данных платёжного приложения
(AID, Application Label) (см. главу \ref{ch:emv}).

Для карт Мир персонализация дополнительно включает установку МРА
(Мир Payment Application) с российской криптографией на базе
ГОСТ (см. главу \ref{ch:mir}).

\subsection*{Authorization host}

Когда авторизационный запрос доходит от карточной сети
до эмитента, его обрабатывает authorization host — система,
принимающая решение: одобрить, отклонить или запросить
дополнительную верификацию. Логика авторизации — цепочка
проверок, выполняемых последовательно:

\begin{enumerate}
  \item \textbf{Статус карты}: не заблокирована ли, не истекла ли.
  \item \textbf{Баланс / лимит}: достаточно ли средств на счёте
        (для дебетовой) или доступного лимита (для кредитной).
  \item \textbf{Лимиты}: не превышен ли дневной, месячный
        или транзакционный лимит.
  \item \textbf{Антифрод}: velocity checks, геолокация, scoring
        модель (см. главу \ref{ch:fraud}).
  \item \textbf{Криптограмма}: для EMV-транзакций —
        верификация ARQC через
        HSM (см. главу \ref{ch:crypto}).
\end{enumerate}

Если все проверки пройдены, host формирует ответ \rc{00}
(Approved) и генерирует ARPC (Authorization Response
Cryptogram), которая передаётся обратно карте для взаимной
аутентификации. Если хотя бы одна проверка не пройдена —
host возвращает соответствующий response code (\rc{51} —
Insufficient Funds, \rc{05} — Do Not Honor, \rc{14} —
Invalid Card Number и так далее), и транзакция отклоняется.

\subsection*{Issuer scripting}

EMV-стандарт предусматривает механизм, позволяющий эмитенту обновлять параметры карты непосредственно через
авторизационный ответ — issuer scripting (Script Processing, команды 71 и 72). Script 71 выполняется до генерации
финальной криптограммы (TC или AAC), Script 72 — после. На практике issuer scripting используется для:

\begin{itemize}
  \item обновления офлайн PIN (смена PIN без визита в банк);
  \item сброса офлайн-счётчиков (после успешной онлайн-авторизации);
  \item блокировки приложения на карте (если эмитент решил
        заблокировать карту, но она уже вставлена в терминал);
  \item обновления параметров риск-менеджмента (лимиты
        офлайн-транзакций).
\end{itemize}

Issuer scripting — мощный инструмент, но его применение ограничено: не все терминалы корректно обрабатывают скрипты, а
при бесконтактной транзакции время на выполнение скриптов ограничено, и сложные обновления могут не успеть завершиться
(см. главу \ref{ch:contactless}).

\subsection*{ACS: Access Control Server}

ACS — компонент эмитента, отвечающий за 3-D Secure аутентификацию (см. главу \ref{ch:3ds}). Когда мерчант инициирует
3DS-проверку, запрос приходит в ACS эмитента, который принимает решение: пропустить транзакцию frictionless (без
участия держателя) или отправить challenge (запрос OTP, биометрия, подтверждение в мобильном приложении).

Решение ACS основано на Risk-Based Authentication (RBA):
анализ устройства, IP-адреса, истории транзакций, суммы,
мерчанта. Качество RBA напрямую влияет на конверсию: слишком
агрессивный ACS отправляет на challenge транзакции, которые
можно было бы пропустить frictionless, теряя покупателей;
слишком мягкий — пропускает фрод, увеличивая убытки
эмитента.

\subsection*{Управление токенами}

Эмитент участвует в жизненном цикле сетевых токенов
(DPAN) (см. главу \ref{ch:tokenization}): получает запрос
на tokenization от TSP (\scheme{Visa} Token Service, MDES,
TSP НСПК), принимает решение об одобрении (на основе
ID\&V — identification and verification держателя),
и управляет состоянием токена — активация, приостановка,
удаление. Когда держатель блокирует карту, эмитент обязан
приостановить все связанные с ней токены; при перевыпуске
карты — обновить привязку токенов к новому PAN.

\subsection*{Dispute management}

Эмитент — первая линия в процессе оспаривания транзакций:
именно к нему обращается держатель карты с жалобой
(\enquote{я не совершал эту покупку}, \enquote{товар
  не получен}). Эмитент проверяет обоснованность претензии
и инициирует chargeback через карточную сеть, направляя его
эквайеру (см. главу \ref{ch:disputes}).

На стороне эмитента dispute management включает: приём
и классификацию претензий, проверку оснований для chargeback
по правилам сети (reason codes, сроки, обязательные
условия), формирование chargeback-сообщения, обработку
representment от эквайера (если эквайер оспорил chargeback,
предоставив compelling evidence) и принятие решения
о дальнейших действиях — принять representment или перейти
к pre-arbitration.

Итак, на стороне эмитента карточная транзакция сводится
к управлению жизненным циклом карты, принятию авторизационного
риска и поддержанию доверенной аутентификации держателя. Теперь
можно зеркально посмотреть на сторону эквайера, где центральной
становится уже не карта, а мерчант и качество его приёма.

\section{Эквайер}
\label{sec:ia-acquirer}

Эквайер (acquirer) — банк, обеспечивающий мерчанту приём
карточных платежей: он предоставляет терминалы (или подключение
к платёжному шлюзу для e-commerce), передаёт авторизационные
запросы в карточную сеть, принимает средства от эмитентов
через settlement и перечисляет их мерчанту за вычетом комиссий.
Если эмитент отвечает на вопрос \enquote{платить ли?},
то эквайер — на вопрос \enquote{как принять?}.

\subsection*{Terminal Management System}

\begin{kaobox}[title={Определение: TMS}, colback=blue!4, colbacktitle=blue!15]
  Terminal Management System (TMS) — система удалённого
  управления парком POS-терминалов: загрузка программного
  обеспечения и обновлений, инъекция криптографических ключей
  (terminal master key, working keys), настройка параметров
  (MCC, Terminal ID, supported card brands, EMV-конфигурация),
  мониторинг состояния и диагностика неисправностей.
\end{kaobox}

Для крупного эквайера TMS управляет десятками и сотнями тысяч
терминалов, и ключевая задача — обеспечить, чтобы все они
работали с актуальной конфигурацией: обновлённые BIN-таблицы,
свежие ключи, корректные EMV-параметры. Удалённая инъекция
ключей (Remote Key Injection, RKI) позволяет загружать ключи
без физического доступа к терминалу, что критично
для масштабирования: ручная key injection ceremony
для каждого терминала нецелесообразна при парке в десятки
тысяч устройств.

\subsection*{Merchant onboarding}

Прежде чем мерчант начнёт принимать платежи, эквайер проводит
onboarding — комплекс проверок и настроек, включающий
несколько этапов.

\textbf{KYB (Know Your Business)} — проверка юридического лица:
регистрация, бенефициары, финансовое состояние, репутация.
Эквайер обязан убедиться, что мерчант не занимается запрещённой
деятельностью, не находится в санкционных списках и имеет
легитимный бизнес. Глубина проверки зависит от уровня риска:
интернет-казино проверяют строже, чем продуктовый магазин.

\textbf{MCC-классификация} — присвоение мерчанту кода
категории (Merchant Category Code), четырёхзначного числа,
определяющего тип бизнеса. MCC — не формальность:
от него зависит ставка interchange (некоторые MCC имеют
пониженные или повышенные ставки), правила 3-D Secure
(exemptions для определённых MCC), лимиты chargeback
и даже легальность операции (MCC 7995 — gambling —
запрещён или ограничен во многих юрисдикциях).

\begin{kaobox}[title=На практике]
  Неправильный MCC — одна из самых распространённых ошибок
  при onboarding, и последствия ощутимы: мерчант платит
  неправильный interchange (переплата или недоплата,
  за которую потом оштрафуют), транзакции некорректно
  проходят антифрод-фильтры, а при аудите карточной сети
  неверный MCC может привести к штрафу и требованию
  рекатегоризации. Уделите MCC-классификации столько же
  внимания, сколько проверке юридических документов.
\end{kaobox}

\textbf{Underwriting} — оценка финансовых рисков: какой объём
транзакций ожидается, каков средний чек, какова вероятность
chargebacks, какой кредитный лимит установить (максимальная
сумма, которую эквайер готов принять от мерчанта до settlement,
неся кредитный риск). Underwriting определяет, на каких
условиях мерчант будет работать: ставка комиссии, holdback
(процент от settlement, удерживаемый как резерв на случай
chargebacks), rolling reserve, payout schedule.

\subsection*{Merchant monitoring}

Onboarding — не конец проверки, а начало. Эквайер обязан непрерывно мониторить мерчанта после подключения:

\begin{itemize}
  \item \textbf{Chargeback ratio} — отношение числа chargebacks
        к числу транзакций. Карточные сети устанавливают пороги
        (VDMP/VFMP для \scheme{Visa}, ECM/EFM
        для \scheme{Mastercard}) (см. главу \ref{ch:disputes}),
        превышение которых влечёт штрафы, повышенный мониторинг
        и в крайнем случае — отключение мерчанта.
  \item \textbf{Fraud rate} — доля мошеннических транзакций.
  \item \textbf{Refund ratio} — аномально высокий процент
        возвратов может сигнализировать о проблемах с товаром
        или о мошеннической схеме (triangulation fraud).
  \item \textbf{Velocity anomalies} — резкий рост объёма
        транзакций, изменение среднего чека, появление
        нехарактерных BIN-диапазонов.
\end{itemize}

\begin{kaobox}[title=Важно, colback=red!5, colbacktitle=red!20]
  Ответственность за мерчанта лежит на эквайере: если мерчант
  оказался мошенником или нарушил правила карточной сети,
  штрафы и финансовые потери несёт эквайер, а не карточная сеть
  и не эмитент. Именно поэтому merchant monitoring — не
  формальность, а критическая функция: пропущенный bust-out
  (мерчант, накопивший объём и внезапно исчезнувший с деньгами)
  может стоить эквайеру миллионов.
\end{kaobox}

\section{Payment Facilitator и sub-merchant модель}
\label{sec:ia-payfac}

Классическая модель эквайринга предполагает, что каждый мерчант
заключает договор с банком-эквайером, получает собственный
MID (Merchant ID), проходит KYB и underwriting. Для крупного
ритейлера это приемлемо, но для онлайн-платформы, которой
нужно подключить тысячи мелких продавцов (маркетплейс,
платформа для фрилансеров, SaaS с подписками), индивидуальный
onboarding каждого продавца через банк — непозволительно
медленный процесс, занимающий дни или недели.

\begin{kaobox}[title={Определение: Payment Facilitator}, colback=blue!4, colbacktitle=blue!15]
  Payment Facilitator (PayFac) — организация, зарегистрированная
  в карточной сети как мерчант-агрегатор, процессящая транзакции
  sub-merchants (подмерчантов) под своим собственным MID.
  PayFac берёт на себя onboarding sub-merchants, распределение
  средств и ответственность за compliance.
\end{kaobox}

В модели PayFac платформа регистрируется в карточной сети
(\scheme{Visa}, \scheme{Mastercard}) как Payment Facilitator,
получает master MID и далее самостоятельно подключает
sub-merchants — каждому присваивается sub-merchant ID,
но все транзакции идут через master MID платформы. Это позволяет
подключить нового продавца за минуты, а не за недели: PayFac
проводит упрощённый KYB, присваивает MCC и начинает принимать
платежи, не дожидаясь банковской бюрократии.

Однако за скорость подключения PayFac платит повышенной
ответственностью: именно PayFac несёт liability
за своих sub-merchants перед карточной сетью и эквайером.
Если sub-merchant оказался мошенником, штрафы ложатся
на PayFac. Если chargeback ratio sub-merchant превысил
порог — отвечает PayFac. Это означает, что PayFac
должен выстроить собственные процессы KYB, merchant
monitoring и dispute management, по строгости не уступающие
банковским.

\begin{kaobox}[title=Важно, colback=red!5, colbacktitle=red!20]
  PayFac — не то же самое, что ISO (Independent Sales
  Organization). ISO — это торговый посредник, который
  привлекает мерчантов для банка-эквайера, но \emph{не} участвует
  в потоке средств и \emph{не} несёт финансовую ответственность
  за мерчанта. PayFac, напротив, стоит в цепочке движения денег:
  settlement приходит на счёт PayFac, и именно PayFac распределяет
  средства между sub-merchants. Это принципиальное различие
  определяет и уровень ответственности, и регуляторные
  требования.
\end{kaobox}

\subsection*{Когда выбирать PayFac}

Модель PayFac оправдана, когда платформе нужно подключить
большое число мелких мерчантов с минимальным friction: маркетплейсы,
платформы для аренды, crowdfunding, SaaS с оплатой по подписке.
Преимущества — скорость onboarding, контроль над пользовательским
опытом, возможность удерживать комиссию и управлять payouts.

Недостатки PayFac-модели:
\begin{itemize}
  \item повышенная ответственность за fraud и chargebacks субмерчантов;
  \item необходимость регистрации в карточных сетях;
  \item инвестиции в собственные KYB/AML-процессы;
  \item в некоторых юрисдикциях — необходимость лицензии
        (в России — статуса платёжного агента или оператора
        по \law{161-ФЗ}).
\end{itemize}

\begin{kaobox}[title=На практике]
  Если вы строите платформу и выбираете между PayFac-моделью
  и прямым эквайрингом для каждого мерчанта, начните
  с вопроса: сколько мерчантов вы планируете подключить
  в первый год? Если десятки — прямой эквайринг проще
  и дешевле. Если сотни и тысячи — PayFac окупится
  за счёт скорости onboarding, но потребует инвестиций
  в compliance-инфраструктуру. Промежуточный вариант —
  использовать готовый PayFac-as-a-Service (Stripe Connect,
  Adyen for Platforms), который берёт на себя регистрацию
  и compliance в обмен на комиссию.
\end{kaobox}

\section*{Итоги главы}

\begin{itemize}
  \item Эмитент управляет жизненным циклом карты через CMS
        (персонализация, статусы, лимиты) и принимает решение
        об авторизации через authorization host (баланс → лимиты
        → антифрод → криптограмма → response code).
  \item Issuer scripting (Script 71/72) позволяет обновлять параметры карты через авторизационный ответ; ACS обеспечивает 3-D
        Secure аутентификацию с решением frictionless/challenge на основе RBA.
  \item Эквайер обеспечивает приём платежей: TMS управляет
        парком терминалов, onboarding включает KYB,
        MCC-классификацию и underwriting, а merchant monitoring
        отслеживает chargeback ratio, fraud rate и velocity
        anomalies.
  \item Payment Facilitator (PayFac) — агрегатор, процессящий sub-merchants под своим MID; даёт скорость onboarding (минуты
        вместо недель), но несёт полную liability за sub-merchants; не путать с ISO, который не участвует в потоке средств.
  \item Выбор между прямым эквайрингом и PayFac определяется
        числом мерчантов и готовностью инвестировать
        в compliance; промежуточный вариант — PayFac-as-a-Service.
\end{itemize}
